\section{Introduction}

% Measuring individual differences within and across species is fundamental to many areas of biology~\parencite{nelson1970outline,wagner1989biological}.
% While such comparisons have yielded significant insights in neuroscience—particularly in surveys of gross-level neuroanatomy across species~\parencite{butler2005comparative} and in studies of small, well-characterized central pattern generator circuits~\parencite{Marder2006}—it remains 
A major challenge in neuroscience is to quantitatively compare the complex, high-dimensional activity of large collections of neural systems. This challenge sits at the center of an ongoing debate about how neural codes are organized, a question that recurs at every scale of description. At the level of single neurons, responses often appear disorganized: neurons typically exhibit mixed selectivity, responding to seemingly arbitrary combinations of variables~\parencite{rigotti2013,tye2024,ostojic2024,posani2026}. At the level of brain regions, a classical view holds that function is localized, with distinct areas playing specialized roles arranged along processing hierarchies (\cite{petersen2024principles}). Recent high-throughput recordings have complicated this picture, revealing that behaviroal variables can be decoded from almost anywhere in the brain~\parencite{steinmetz2019distributed,stringer2019spontaneous,musall2019, findling2025brain,angelaki2025brain,rosen2026distributed,schartner2025, bondy2025,posani2026}, a result interpreted as evidence that population codes are broadly distributed and largely undifferentiated across regions (\textcite{hayden2026rethinking,fine2021whole,Lloreda2025}). Notably, these conclusions often rest on summaries of single-neuron tuning pooled
across many subjects, which presupposes that variability from one animal to the next is idiosyncratic noise rather than structured signal. 

Part of these debates could be resolved by moving from single-neuron tuning and single-variable decodability to a fully geometrical picture, i.e. how the full set of behavioral variables is jointly arranged and how these arragentments differ within and across individual. Indeed, recent work has argued that this arrangement, often called neural geometry, directly reflects a system's computation (\cite{courellis2024,nogueira2026, fascianelli2024neural, pagan2025}).

\begin{figure}
    \centering
    \includegraphics[width=0.99\textwidth]{content/figures/netrep-journal-figure-1.pdf}
    \caption{\textbf{A geometric framework for comparing neural populations at scale.} \textbf{a-c.} Shape metrics can be leveraged to compare representations across animals (a), regions (b) or different time periods from the same animal (c). \textbf{d.} Example population activity measured across $M$ conditions and $N$ neurons in an experiment parameterized by a periodic 1D variable, $\theta$ (e.g. head direction). Smooth tuning curves, \textit{left}, can be viewed as a smooth manifold, \textit{right}, in $N$-dimensional firing rate space. Color code denotes values of $\theta$. \textbf{e.} The goal of our framework is to compute a distance function $d(X_i,X_j)$, which we define as the dissimilarity between two manifolds $X_i, X_j$ after aliment. \textbf{f.} The set of possible manifolds and the distance function together define a \textit{metric space}. \textbf{g.} The proprieties of these metric spaces allow us to leverage  existing methods, such as clustering and regression with rigorous theoretical guarantees.}
    \label{fig:schematic}
\end{figure}

However, comparing that geometry across systems has itself proven methodologically contentious. Over the past two decades numerous measures have been proposed, beginning with Representational Similarity Analysis (RSA; \cite{Kriegeskorte2008_frontiers}) and its many variants \parencite[see][for a review]{kriegeskorte2019peeling}, alongside Canonical Correlations Analysis \parencite[CCA;][]{Raghu2017,gallego2020long} and hyperalignment \parencite{haxby2011common,haxby2020hyperalignment}, as well as regression- and classification-based comparisons of representational geometry \parencite{bernardi,conwell2022can,fascianelli2024neural,muzellec2026reverse}. Crucially, scaling up these analyses to compare large collections of neural systems---across subjects, sessions, and brain regions (\textit{Fig.}\ref{fig:schematic}a-c)---remains a challenge.
For instance, one might ask whether the recordings fall into clusters, e.g. small groups of subjects, or of brain regions, whose population activity sits closer to one another's than to the rest. An affirmative answer would establish that variability across recordings is structured rather than purely idiosyncratic. Similar questions have been asked at the level of single neuron tuning proprieties (\cite{hirokawa2019,posani2026}), but only rarely for entire recordings, and even then for small numbers of subjects \parencite{safaie2023preserved,Chen2021,gallego2020long}. One might also ask whether the arrangement of these systems are behaviorally relevant. Such question could be addressed by testing whether a subject's behavioral profile can be inferred from other subjects with similar neural codes. Answering either types question presuppose that the systems occupy a common space, each recording holding a position within it and distances across systems reflecting their similarity.


Here, we leverage a principled approach to geometrically compare activity measurements among multiple neural populations, such as many different brain regions or animal subjects, which is broadly referred to as \emph{shape distance} (\cite{Williams2021}). Under this framework, an entire collection of recordings becomes a cloud of points (\textit{Fig.}\ref{fig:schematic}f), each neural system holding a position relative to every other, over which clustering and regression can be performed rigorously (\textit{Fig.}\ref{fig:schematic}g). The mathematical foundations of the shape metrics framework were laid out in a series of recent conference papers \parencite{Williams2021,pospisil2023estimating,harvey2024duality,khosla2024soft}.
In this work, we applied this framework to show that neural population geometry is structured within and across individuals.

Applying this framework across dozens of brain regions in mice and non-human primates, we recover a hierarchical organization of brain regions in a fully unsupervised manner, despite that individual task variables are decodable from nearly everywhere in these brains. In the mouse head direction system, we uncover systematic animal-to-animal differences in head direction coding manifolds, trace these differences to tuning width via regression in shape space, and show that all subjects nonetheless share a common ring-shaped manifold, enabling decoding of head direction in a held-out animal from other animals' data alone \parencite{safaie2023preserved}. Further, in a large cohort of mice performing a decision-making task, we show that individual differences in behavior are mirrored by individual differences in neural population geometry in most considered regions, allowing one animal's behavior to be predicted from the neural data of its closest neighbors. Finally, we find that genetically defined differences between animals, invisible at the single-neuron level~\parencite{noel2025}, are reflected in population geometry.

Together, these results show that neural population geometry is structured across both regions and individuals, organized by anatomy across brain areas and by behavior and genotype across individuals.